\documentclass{article}

\usepackage[pattern=dot, 
patternsize=0.2in, 
textarea,
geometry={
	a4paper,
	left=1.5cm,
	right=1.5cm,
	top=15mm,
	bottom=15mm,
	headsep=10pt,
	footskip=25pt,
}
]{gridpapers}
\usepackage{fancyhdr}
\pagestyle{fancy}
\usepackage[normalem]{ulem}
\usepackage{cancel}
\usepackage[dvipsnames]{xcolor}
\usepackage{parskip}
\usepackage{caption, subcaption}
\usepackage[Spanish,es-tabla,es-noshorthands]{babel}
\usepackage{newtxtext}
\usepackage[varvw]{newtxmath}
\usepackage{booktabs, colortbl, bigstrut, multirow, multicol}
\usepackage{url}
\usepackage{pgfplots}
\usepackage{graphicx}
\usepackage{amsmath}
\usepackage{tabularx}
\usepackage{floatrow}
\usepackage{float}
\usepackage{cancel}
\usepackage[spanish]{cleveref}
\usepackage{tikz}
\usetikzlibrary{arrows.meta,shapes.geometric,positioning,calc}
\crefname{table}{tabla}{tablas}
\usepackage{wrapfig}
\usepackage{adjustbox}
\usepackage{fancybox}
\usepackage[T1]{fontenc}
\renewcommand*\familydefault{\sfdefault}


% Commands
\renewcommand{\headrulewidth}{0.4pt}
\renewcommand{\footrulewidth}{0.4pt}
\let\HeadRule\headrule
\let\FootRule\footrule
\renewcommand\headrule{\color{gray}\HeadRule}
\renewcommand\footrule{\textcolor{gray}{\FootRule}}

\newcommand\hcancel[2][black]{\setbox0=\hbox{$#2$}%strikeout
\rlap{\raisebox{.45\ht0}{\textcolor{#1}{\rule{\wd0}{1pt}}}}#2} 

\newsavebox\MBox
\newcommand\Cline[2][red]{{\sbox\MBox{$#2$}%
  \rlap{\usebox\MBox}\color{#1}\rule[-1.2\dp\MBox]{\wd\MBox}{0.5pt}}}

\renewcommand\theequation{\Alph{equation}}

%\renewcommand\thesubsection{\alph{subsection}}

\def\mydate{\leavevmode\hbox{\twodigits\day/\twodigits\month/\the\year}}
\def\twodigits#1{\ifnum#1<10 0\fi\the#1}

\newcommand\runinsec{\@startsection {section}{1}{\z@}%
                                   {3.5ex \@plus 1ex \@minus .2ex}%
                                   {-1em}%
                                   {\normalfont\Large\bfseries}}
\newcommand\aftersec[1]{\begingroup\normalfont\footnotesize #1\endgroup%
                        \par\nobreak\vspace*{2.3ex}\noindent\ignorespaces}


% Header and footer
\lhead{\textcolor{gray}{MM1 -- \thetitle}}
\rhead{\textcolor{gray}{Date: \mydate \hspace{1cm} Page No.: \thepage}}
\cfoot{\textcolor{gray}{\thefooter}}

\input{colorbox.txt}

% Document-------------------------------------

\newcommand{\thetitle}{Resolución de Ecuaciones Diferenciales}
\newcommand{\thefooter}{}

\begin{document}
\ \\ \vspace{-0.6cm} 
{\Huge\textsc{\thetitle}} \vspace{0.3cm}
\input{tikz.txt}

\begin{wrapfigure}[7]{r}{0.7\textwidth}
\vspace{-1.66cm}
\begin{center}
\begin{tikzpicture}[node distance=1cm]


\node [txt] (ex) {$ \left(
	\begin{array}{l}
		\text{(Exacta)}\\ \textcolor{gray}{ \sim\text{\small fac. int.}}
	\end{array}\right)$
};
\node [rect, above= of ex] (lin) {Lineal};
\node[rect,right=of lin, xshift=50px](cla){Clairaut};
\node[rect,above=of lin,xshift=-2.6cm](ber){Bernoulli};
\node[rect,above=of lin](ric){Ricatti \\ {\small (+ sol. conocida)}};
\node[rect,above=of lin,xshift=2.6cm](hom){Homogénea\\ (cocientes)};
\node[ellip,below=of cla](sol){\textbf{\textsc{Solución}}\\ \textbf{\textsc{General}}} node[below=of sol, yshift=1cm]{(+ Condiciones iniciales)};
\node [txt, left=of ber, xshift=1cm] (ini) {\textbf{\underline{¿Tipo?:}}};

\draw[arrow] (lin)--(ex) ;
\draw[arrow] (ber)|-(lin) node[midway,right,yshift=30px] {$v=y^{1-n}$};
\draw[arrow] (ric)--(lin) node[midway,right,yshift=2px] {$y=y_p+\frac{1}{v}$};
\draw[arrow] (hom)|-(lin) node[midway,right,yshift=30px] {$v=\frac{y}{x}$};
\draw[arrow] (cla)--(sol) node[midway,right] {$y'=p$};
\draw[arrow] (ex)--(sol) node[midway,above] {$\Phi$} node[midway,below] {\textcolor{gray}{(¿separable?)}};
\draw[tarrow] (lin)--(sol) node[midway, above, sloped]{$\mu (x)$};

\end{tikzpicture}
\end{center}
\end{wrapfigure}

\vspace{0.8cm}
\section{\textsc{Primer orden:}}

Para resolver una \textit{EDO} debemos poner una $y(x)$ en función de su \textbf{variable independiente} y \textbf{constantes}. Para ello se emplean \underline{cambios de variable}, que luego se derivan y sustituyen en la ecuación principal. Existen varios definidos:
\vspace{10pt}

\subsection{\underline{Lineal} \aftersec{\colorbox{CadetBlue!20}{$\; y'+P(x)y=Q(x)$}}}
\vspace{-0.7cm}
Se \textbf{multiplica} toda la ecuación por el factor integrante \colorbox{pink}{$\mu(x)=e^{\int{Pdx}}$} y se resuelve.

\subsection{\underline{Exacta} \aftersec{(cumplen $\frac{\partial P}{\partial y}=\frac{\partial Q}{\partial x}$)}}
\vspace{-0.6cm}
\subsubsection{Separables}

Si se pueden \textbf{separar} las variables a cada lado, se puede resolver simplemente \textbf{integrando ambos lados} para obtener la función $y(x)$ .

\subsubsection{No separables}

Hay que usar la propiedad de que, al ser exacta, con forma \underline{$P(x,y)dx+Q(x,y)dy=0$} , existe una función cuya \textbf{derivada} cumple $dF=Pdx+Qdy$, siendo $P$ y $N$ ahora constantes.

\begin{colorboxdef}{Exactas no separables}{gray}
	Por definición tenemos que $\boxed{P={\partial F}/{\partial x}}$ y $\boxed{Q={\partial F}/{\partial y}}$
	\begin{itemize}
	\item Despejamos $\partial F$ de cualquiera e integramos todo para hallar $F$\\
	$\hookrightarrow$ La \textbf{constante de int.} será $g(x)$ o $g(y)$
	\item Para hallar la $g$ que nos falta, emplear la otra expresión no usada \textbf{derivando} $\frac{\partial F}{\partial x\, \text{ó} \, \partial y}$ e igualar a la $P\, \text{ó}\, Q$ correspondiente
	\item \textbf{Susituir} la $g$ hallada en la $F$ obtenida integrando inicialmente $\to$ \underline{\textsc{solución}}
\end{itemize}
\end{colorboxdef}

\subsection{\underline{Bernoulli} \aftersec{\colorbox{CadetBlue!20}{$\; y'=P(x)y=Q(x)y^n$}}}
\vspace{-0.7cm}
Realizamos un cambio de variable \colorbox{pink}{$v=y^{1-n}$} y \textbf{dividimos o multiplicamos} toda la ecuación para que termine como una \textbf{lineal} en función de $v$ y $v'\, \to$ \textit{resolvemos la Lineal}.

\subsection{\underline{Homogénea} (Cocientes) \aftersec{\colorbox{CadetBlue!20}{$\; y'=f(y/x)\;$} }}
\vspace{-0.7cm}
Se hace un cambio de variable \colorbox{pink}{$v=\frac{y}{x}$}
\vspace{-0.3cm}

\subsection{\underline{Ricatti} \aftersec{\colorbox{CadetBlue!20}{$\; y'=P(x)y^2+Q(x)y+R(x)$} }}
\vspace{-0.7cm}
Si se conoce una \textbf{solución particular} ($y_p$), se puede hacer el cambio \colorbox{pink}{$y=y_p+\frac{1}{v}$}
\vspace{-0.3cm}

\subsection{\underline{Clairaut} \aftersec{\colorbox{CadetBlue!20}{$\; y=xy'+f(y')$}}}
\vspace{-0.7cm}

Definimos \colorbox{pink}{$y'=p$} y \textbf{derivamos todo} en función de $x$.\\
$\hookrightarrow$ Acabaremos con una función $A(B)=0$ , que podemos parametrizar: $ \left\{
	\begin{array}{l}
		A=0\\
		B=0
	\end{array}
\right. $\\
$\hookrightarrow$ Una solución nos dará la \underline{general} y con la otra es posible desarrollar las paramétricas y hallar una \underline{singular}.


\begin{figure}[h!]
	\centering
	\begin{minipage}{.69\textwidth}
		\section{\textsc{Orden superior:}}
		\subsection{\underline{Solución General} \hspace{0.5cm} \aftersec{$(Ay''+By'+cY'=g(x))$}}
\vspace{-0.72cm}
Es la suma de la solución particular y la de la homogénea asociada: $\boxed{y(x)=y_h+y_p}$ . 
	\end{minipage}
	\begin{minipage}{.3\textwidth}
\raggedright
\includegraphics[width=\columnwidth]{img-hom/file.pdf}
	\end{minipage}
	\end{figure}
\vspace{-0.8cm}


\subsubsection{Homogénea -- \textsc{\normalsize coeficientes constantes}}

Las homogéneas tienen $g(x)=0$, siendo $g(x)$ una función dependiente de $x$ a la derecha de la ecuación.

\begin{colorboxdef}{Solucion \textbf{Homogénea} -- coeficientes \textbf{Constantes}:}{Plum}
	Tienen la forma: $Ay''+By'+cY'=0$, y se resuelven con su ecuación característica:\\
	\begin{center}
		 \underline{$y''=r^2$}\; \underline{$y'=r$}\; \underline{$y=1$}
	\end{center}
\vspace{10pt}

	La solución \textbf{homogénea} tendrá forma distinta dependiendo de las raíces halladas:
	\begin{align*}
		\left\{
		\begin{array}{lcr}
			\text{\textbf{Raíces distintas:} }\; &y=C_1e^{r_1x}+C_2e^{r_2x}\\
			\text{\textbf{Raiz doble:} }\; &y=C_1e^{rx}+C_2\mathcolor{red}{x}e^{rx}& \text{\textcolor{gray}{(por cada sol. multiplic)}}\\
			\text{\textbf{Raíces imag.:} }\; &y=e^{\alpha x}(C_1\cos{(\beta x)}+C_2\sin{(\beta x)})& \hspace{20pt} \mathcolor{gray}{(r=\alpha\pm\beta i)}\\
		\end{array}
		\right.
	\end{align*}
\end{colorboxdef}

\subsubsection{Solución particular}	

Tras obtener la solución de la homogénea asociada, se halla la solución particular:
\begin{colorboxdef}{Solución \textbf{Particular} -- \textit{Coeficientes indeterminados}}{Plum}
	Se prueba una forma de $y_p$ según el caso:
\begin{gather*}
	\left\{
	\begin{array}{l}
		y_p=Ax^2+Bx+C\\
		y_p=Ae^{2x}\\
		y_p=A\cos x +B\sin x
	\end{array}
	\right.
\end{gather*}
\end{colorboxdef}

Si la ecuación tiene varias características, se emplean varios tipos de solución y se \textbf{multiplican} entre ellos. Existe un método universal si el caso no está claro:

\begin{colorboxdef}{Método Universal \textbf{Particular} -- \textit{Variación de parámetros}}{lightgray}
La $y_p$ tendrá una forma genérica \underline{$y_p=u_1(x)y_1(x)+u_2(x)y_2(x)$} . Se imponen las siguientes condiciones:
	\begin{gather*}
	\left\{
	\begin{array}{l}
		u_1'y_1+u_2'y_2=0\\
		u_1'y_1'+u_2'y_2'=g(x)
	\end{array}
	\right. \; \text{ y se \textbf{resuelve el sistema.}}
	\end{gather*}
\end{colorboxdef}

\subsection{\underline{Coeficientes Variables $\to$ Constantes}}

Cuando los coeficientes \textbf{no son constantes}, es necesario manipular la ecuación y/o realizar \underline{cambios de variable} para conseguir que lo sean.

Por lo general, escogemos un cambio que sepamos que, una vez derivado las veces que sea, va a eliminar un factor de cada coeficiente para acabar simplificando la ecuación. Se emplea la propiedad $\boxed{\frac{dy}{dx}=\frac{dy}{dt}\frac{dt}{dx}}$ , siendo $t$ la variable del cambio. 

$\hookrightarrow$ \textcolor{gray}{Suele ser útil el \underline{$Ln$}, ya que su derivada es $\frac{1}{t'}$, que ayuda a eliminar factores. }

$\hookrightarrow$ \textcolor{gray}{Cuidado con la \underline{regla de la cadena}, ya que: $\; \frac{\partial^2 y}{\partial x^2}=\frac{\partial }{\partial x}\left(\frac{\partial y}{\partial x}\right) = \frac{\partial }{\partial x}
	\tikz[baseline]{
		\node(dydt){$\displaystyle \left(\frac{\partial y}{\partial t}\frac{\partial t}{\partial x}\right)$};
		\node[below=of dydt,yshift=20pt,xshift=85pt,red](dydtAn){$\displaystyle \frac{\partial}{\partial x}\left(\frac{\partial y}{\partial t}\right)=\frac{\partial}{\partial t}\left(\frac{\partial y}{\partial t}\right)\frac{\partial t}{\partial x}=\frac{\partial^2 y}{\partial t^2}\frac{\partial t}{\partial x}$};
		\draw [->,red] (dydt.south) |- (dydtAn.west) node[midway, left]{\text{\textbf{(r.c. 2)}}};
		\node[right=of dydt,xshift=-30pt](dydtN){$\displaystyle \xrightarrow[(\text{r. cadena})]{} \left( \Cline[red]{\frac{\partial^2 y}{\partial t^2}\frac{\partial t}{\partial x}} \right)\frac{\partial t}{\partial x} +\frac{\partial y}{\partial t}\frac{\partial^2 t}{\partial x^2}$};
	}
	$}


\subsubsection{Cauchy-Euler \aftersec{\colorbox{CadetBlue!20}{$\; ax^2y''+bxy'+cy=0$}}}
\vspace{-0.7cm}
Las ecuaciones de Cauchy tienen un cambio de variable definido (al igual que \textit{Bernoulli}, \textit{Ricatti}\dots), que puede ayudar a simplificar el caso.

	\begin{itemize}
	\item Hacemos un cambio de variable \colorbox{pink}{$y=x^r$} y derivamos como sea necesario para \textbf{sustituir}
	\item Agrupamos, manipulamos\dots para \textbf{resolver} la ecuación y hallar las \textbf{raíces} ($r$)
	\item Aplicamos soluciones según las raíces \[ \boxed{\left\{
	\begin{array}{lc}
		\text{\textbf{Raíces distintas:} }\; &y=C_1x^{r_1}+C_2x^{r_2}\\
			\text{\textbf{Raiz doble:} }\; &y=C_1x^{r}+C_2x^{r}\mathcolor{red}{\ln x}\\
			\text{\textbf{Raíces imag.:} }\; &y=x^{\alpha}(C_1\cos{(\beta \ln x)}+C_2\sin{(\beta \ln x)})
	\end{array}
	\right. }
	\]
	\end{itemize}


\subsection{\underline{Reducción de Orden} \aftersec{(solución $y_1$ conocida)}}
\vspace{-20pt}
\begin{itemize}
	\item El cambio de variable $\boxed{y_2=y_1 v(x)}$ reduce el orden en \textbf{1 ud.}
	\item Derivamos $y'$, $y''$\dots y \textbf{sustituimos} en la principal
	\item Realizamos el cambio de variable $\boxed{u=v'}$, resultando en una ecuación de $1^{er}$ orden\\
	$\hookrightarrow$ \textbf{Resolvemos} la ecuación\\
	$\hookrightarrow$ \textbf{Deshacemos cambio} integrando
\end{itemize}

\section{\textsc{Herramientas}}

\begin{itemize}
	\item \textbf{Comprobar si es solución particular:} derivamos las veces necesarias para sustituir en la ecuación principal y ver si se cumple la igualdad.
	\item \textbf{Condiciones iniciales:} Para hallar las constantes de integración, aplicamos las condiciones iniciales dadas, si hay.
	\item \textbf{\textsc{Wronskiano}:} si \underline{$W(x_o) \neq 0$} , se sabe que dos soluciones son \textbf{independientes}. \[W(y_1,y_2)(x)=  \left|
	\begin{array}{cc}
	y_1&y_2\\y_1'&y_2'
	\end{array} \right| \]
\end{itemize}




\end{document}